\documentclass[12pt]{article}
\usepackage[marginpar=0pt]{geometry}
\geometry{verbose,tmargin=2.5cm,bmargin=3cm,lmargin=1.5cm,rmargin=1.5cm,footskip=1cm}
\usepackage{amsfonts, amsmath, amssymb, amsthm,graphicx}
\usepackage{physics}
\usepackage{xcolor}
\usepackage{hyperref}
\usepackage{tabularx}
\usepackage{authblk}
\usepackage{dsfont}

\usepackage[hang]{footmisc}
\setlength{\footnotemargin}{1em} % indent after the footnote number
\renewcommand{\baselinestretch}{1.1}  % between the lines

% custom total width of the footnote block
\renewcommand{\footnotelayout}{%
  \setlength{\parindent}{0pt}% no extra indent
  \setlength{\leftskip}{2em}% shift whole block to the right
  \setlength{\rightskip}{3em}% optional right margin
}
\usepackage{cite}
\usepackage{caption}
%\captionsetup[figure]{width=0.85\linewidth}
%\captionsetup[figure]{font=footnotesize} 

\newtheorem{proposition}{Proposition}
\newtheorem{definition}{Definition}
\newtheorem{theorem}{Theorem}
\newtheorem{lemma}{Lemma}
\newtheorem{corollary}{Corollary}

\newcommand{\id}{\mathds{I}}

\newcommand{\cA}{{\mathcal A}}
\newcommand{\cB}{{\mathcal B}}
\newcommand{\cC}{{\mathcal C}}
\newcommand{\cD}{{\mathcal D}}
\newcommand{\cE}{{\mathcal E}}
\newcommand{\cF}{{\mathcal F}}
\newcommand{\cG}{{\mathcal G}}
\newcommand{\cH}{{\mathcal H}}
\newcommand{\cI}{{\mathcal I}}
\newcommand{\cJ}{{\mathcal J}}
\newcommand{\cK}{{\mathcal K}}
\newcommand{\cL}{{\mathcal L}}
\newcommand{\cM}{{\mathcal M}}
\newcommand{\cN}{{\mathcal N}}
\newcommand{\cO}{{\mathcal O}}
\newcommand{\cP}{{\mathcal P}}
\newcommand{\cQ}{{\mathcal Q}}
\newcommand{\cR}{{\mathcal R}}
\newcommand{\cS}{{\mathcal S}}
\newcommand{\cT}{{\mathcal T}}
\newcommand{\cU}{{\mathcal U}}
\newcommand{\cV}{{\mathcal V}}
\newcommand{\cW}{{\mathcal W}}
\newcommand{\cX}{{\mathcal X}}
\newcommand{\cY}{{\mathcal Y}}
\newcommand{\cZ}{{\mathcal Z}}

\renewcommand\a{\alpha}
\renewcommand\b{\beta}
\renewcommand\c{\chi}
\newcommand\C{\Chi}
\renewcommand\d{\delta}
\newcommand\D{\Delta}
\newcommand\e{\epsilon}
\newcommand\ve{\varepsilon}
\newcommand\g{\gamma}
\newcommand\G{\Gamma}
\newcommand\h{\eta}
\renewcommand\k{\kappa}
\renewcommand\l{\lambda}
\renewcommand\L{\Lambda}
\newcommand\m{\mu}
\newcommand\n{\nu}
\newcommand\p{\phi}
\renewcommand\P{\Phi}
\newcommand\vp{\varphi}
\newcommand\q{\theta}
\newcommand\Q{\Theta}
\renewcommand\r{\rho}
\newcommand\s{\sigma}
\renewcommand\S{\Sigma}
\renewcommand\t{\tau}
\renewcommand\u{\upsilon}
\newcommand\U{\Upsilon}
\newcommand\w{\omega}
\newcommand\W{\Omega}
\newcommand\x{\xi}
\newcommand\X{\Xi}
\newcommand\y{\psi}
\newcommand\Y{\Psi}
\newcommand\z{\zeta}



\begin{document}
Lagrange density:
\begin{equation}
    \mathcal{L} = \frac{1}{2}\tr\left( \dot{\Phi}^2 -(\nabla\Phi)^2 - m^2 \Phi^2 - 2g^2 \Phi^4\right).
\end{equation}
Squared commutator:
\begin{equation}
    C(t) = \frac{1}{N^4}\sum_{aba'b'}\int d^3\mathbf{x} \ \tr\Big( \sqrt{\rho}\, [\Phi_{ab}(t,\mathbf{x}),\Phi_{a'b'}]\sqrt{\rho}\, [\Phi_{ab}(t,\mathbf{x}),\Phi_{a'b'}]^\dag\Big).
\end{equation}
Retarded propagator $G_R$ and Wightman function $\tilde{G}$:
\begin{align}
\delta_{ab'}\delta_{ba'}G_R(\mathbf{x},t)&= \theta(t) \tr\Big(\rho\,  [\Phi_{ab}(\mathbf{x},t),\Phi_{a'b'}]\Big)\\
\delta_{ab'}\delta_{ba'}\tilde{G}(\mathbf{x},t)&=\tr\Big(\sqrt{\rho}\, \Phi_{ab}(\mathbf{x},t)\sqrt{\rho}\,\Phi_{a'b'}\Big) .
\end{align}
If Fourier space these take the following form in terms of the spectral
function $\r(k) = \r(k^0,|\vec{k}|)$:
\begin{align}
G_R(k) &= i\int \frac{d\omega}{2\pi}\frac{\rho(\omega, |\k|)}{k^0 - \omega + i\epsilon}\\
\tilde{G}(k) &= \frac{\rho(k)}{2\sinh \frac{\beta k^0}{2}}.
\end{align}
For a free field, we have $\rho(\omega,\k) = \frac{\pi}{E_{\k}}[\delta(\omega-E_{\k}) - \delta(\omega + E_{\k})]$, and therefore
\begin{align}
G_R(k) &= \frac{i}{2E_{\k}}\left(\frac{1}{k^0 - E_{\k}+i\epsilon} - \frac{1}{k^0 + E_{\k} + i\epsilon}\right)\label{retardedG}\\
\tilde{G}(k) &= \sum_{s = \pm}\frac{\pi\delta(k^0-sE_\k)}{2E_\k\sinh\frac{\beta E_\k}{2}}.\label{tildeG}
\end{align}

\end{document}
